\begin{abstract}
Conditional Flow Matching (CFM) trains continuous normalizing flows via simulation-free velocity regression, but the single-sample conditional target induces high gradient variance. Three independent techniques address this: minibatch optimal-transport coupling (BatchOT), importance-weighted marginal velocity targets (StableVM), and antithetic temporal sampling with consistency regularization (TPC). We investigate whether these techniques compose additively through two complementary studies. On 2D synthetic distributions, we confirm that OT-type paths produce straighter trajectories and superior low-NFE robustness. On CIFAR-10, we cross the three axes in a $2^3$ factorial ablation: only one of eight cells converged---BatchOT with standard CFM and uniform sampling (FID~18.4). The other seven diverged or oscillated due to dimension-dependent numerical instabilities invisible in 2D. These negative results reveal that StableVM's importance weights scale catastrophically with dimension ($d{=}3072$), TPC's consistency penalty causes gradient explosions, and the three axes are not additive. BatchOT alone is the single most impactful variance-reduction technique.
\end{abstract}
